\providecommand{\weakconv}{\rightharpoonup}

\subsection{Свойства компактных операторов}

\begin{definition}[Предкомпактное множество]
	Пусть \(E\) --- нормированное пространство. Множество
	\[
	M\subset E
	\]
	называется \textbf{предкомпактным}, если его замыкание
	\[
	\overline{M}
	\]
	компактно.
\end{definition}

\begin{definition}[Компактный оператор]
	Пусть \(E_1,E_2\) --- нормированные пространства и
	\[
	A\in\mathcal{L}(E_1,E_2).
	\]
	Оператор \(A\) называется \textbf{компактным}, если образ любого ограниченного множества
	в \(E_1\) является предкомпактным множеством в \(E_2\).
	
	Множество компактных операторов обозначают
	\[
	\mathcal{K}(E_1,E_2).
	\]
	Если \(E_1=E_2=E\), пишут
	\[
	\mathcal{K}(E).
	\]
\end{definition}

\begin{remark}[Предкомпактность и вполне ограниченность]
	В банаховом пространстве предкомпактность эквивалентна вполне ограниченности.
	
	То есть множество \(M\) предкомпактно тогда и только тогда, когда для любого
	\[
	\varepsilon>0
	\]
	существует конечная \(\varepsilon\)-сеть для \(M\).
\end{remark}

\begin{lemma}[Критерий компактности через единичный шар]
	Пусть \(E_1,E_2\) --- нормированные пространства и
	\[
	A\in\mathcal{L}(E_1,E_2).
	\]
	Тогда \(A\) компактен тогда и только тогда, когда множество
	\[
	A\bigl(\overline{B}_1(0)\bigr)
	\]
	предкомпактно в \(E_2\).
\end{lemma}

\begin{proof}
	Докажем обе стороны.
	
	\textbf{1. Если \(A\) компактен, то \(A(\overline{B}_1(0))\) предкомпактно.}
	
	Единичный шар
	\[
	\overline{B}_1(0)
	\]
	ограничен. Поэтому по определению компактного оператора его образ
	\[
	A\bigl(\overline{B}_1(0)\bigr)
	\]
	предкомпактен.
	
	\textbf{2. Если \(A(\overline{B}_1(0))\) предкомпактно, то \(A\) компактен.}
	
	Пусть
	\[
	S\subset E_1
	\]
	--- произвольное ограниченное множество. Тогда существует \(R>0\), такое что
	\[
	S\subset R\overline{B}_1(0).
	\]
	Так как \(A\) линеен,
	\[
	A(S)\subset A\bigl(R\overline{B}_1(0)\bigr)
	=
	R A\bigl(\overline{B}_1(0)\bigr).
	\]
	Если \(A(\overline{B}_1(0))\) предкомпактно, то и его растяжение
	\[
	R A\bigl(\overline{B}_1(0)\bigr)
	\]
	предкомпактно. Следовательно, \(A(S)\) как подмножество предкомпактного множества тоже
	предкомпактно.
	
	Значит, \(A\) компактен.
\end{proof}

\begin{remark}[Последовательностный критерий компактности]
	Оператор
	\[
	A\in\mathcal{L}(E_1,E_2)
	\]
	компактен тогда и только тогда, когда для любой ограниченной последовательности
	\[
	\{x_n\}\subset E_1
	\]
	последовательность
	\[
	\{Ax_n\}\subset E_2
	\]
	имеет сходящуюся подпоследовательность.
\end{remark}

\begin{theorem}[Конечномерный случай]
	Пусть \(E_1,E_2\) --- нормированные пространства и
	\[
	A\in\mathcal{L}(E_1,E_2).
	\]
	Если
	\[
	\dim E_1<\infty
	\]
	или
	\[
	\dim E_2<\infty,
	\]
	то оператор \(A\) компактен.
\end{theorem}

\begin{proof}
	\textbf{1. Если \(\dim E_1<\infty\).}
	
	Тогда образ оператора
	\[
	\operatorname{Im}A
	\]
	тоже конечномерен, потому что
	\[
	\dim\operatorname{Im}A\leqslant \dim E_1.
	\]
	Если \(S\subset E_1\) ограничено, то \(A(S)\) ограничено в конечномерном пространстве
	\(\operatorname{Im}A\). В конечномерном пространстве всякое ограниченное множество
	предкомпактно. Следовательно, \(A(S)\) предкомпактно.
	
	Значит, \(A\) компактен.
	
	\textbf{2. Если \(\dim E_2<\infty\).}
	
	Пусть \(S\subset E_1\) ограничено. Так как \(A\) ограничен,
	\[
	A(S)
	\]
	тоже ограничено. Но \(A(S)\subset E_2\), а \(E_2\) конечномерно. Поэтому \(A(S)\)
	предкомпактно.
	
	Значит, \(A\) компактен.
\end{proof}

\begin{theorem}[Тождественный оператор не компактен в бесконечномерном пространстве]
	Пусть \(E\) --- бесконечномерное нормированное пространство. Тогда
	\[
	I\notin\mathcal{K}(E).
	\]
\end{theorem}

\begin{proof}
	Предположим противное:
	\[
	I\in\mathcal{K}(E).
	\]
	Тогда образ единичного шара под действием \(I\) предкомпактен:
	\[
	I\bigl(\overline{B}_1(0)\bigr)=\overline{B}_1(0).
	\]
	Значит, единичный шар был бы предкомпактен.
	
	Но по теореме Рисса в бесконечномерном нормированном пространстве единичный шар не
	является компактным и даже не является вполне ограниченным.
	
	Получили противоречие. Следовательно,
	\[
	I\notin\mathcal{K}(E).
	\]
\end{proof}

\begin{theorem}[Компактные операторы образуют двусторонний идеал]
	Пусть \(E\) --- нормированное пространство. Тогда
	\[
	\mathcal{K}(E)
	\]
	является двусторонним идеалом в \(\mathcal{L}(E)\).
	
	То есть если
	\[
	A\in\mathcal{K}(E),
	\qquad
	B\in\mathcal{L}(E),
	\]
	то
	\[
	AB\in\mathcal{K}(E),
	\qquad
	BA\in\mathcal{K}(E).
	\]
\end{theorem}

\begin{proof}
	Докажем компактность обоих произведений.
	
	\textbf{1. Компактность \(AB\).}
	
	Пусть
	\[
	S\subset E
	\]
	--- ограниченное множество. Так как \(B\in\mathcal{L}(E)\), оператор \(B\) ограничен. Поэтому
	\[
	B(S)
	\]
	тоже ограничено.
	
	Так как \(A\) компактен, образ ограниченного множества \(B(S)\) под действием \(A\)
	предкомпактен:
	\[
	A(B(S))
	\text{ предкомпактно}.
	\]
	Но
	\[
	A(B(S))=(AB)(S).
	\]
	Значит, \(AB\) компактен.
	
	\textbf{2. Компактность \(BA\).}
	
	Пусть
	\[
	S\subset E
	\]
	--- ограниченное множество. Так как \(A\) компактен,
	\[
	A(S)
	\]
	предкомпактно.
	
	Нужно доказать, что
	\[
	B(A(S))
	\]
	предкомпактно.
	
	\begin{itemize}
		\item \textbf{Через компактное замыкание.}
		
		Так как \(A(S)\) предкомпактно, множество
		\[
		\overline{A(S)}
		\]
		компактно.
		
		Оператор \(B\) ограничен, значит непрерывен. Поэтому образ компакта при \(B\)
		компактен:
		\[
		B\bigl(\overline{A(S)}\bigr)
		\text{ компактен}.
		\]
		
		\item \textbf{Переход к нужному множеству.}
		
		Имеем
		\[
		B(A(S))\subset B\bigl(\overline{A(S)}\bigr).
		\]
		Следовательно, \(B(A(S))\) является подмножеством компактного множества, значит
		предкомпактно.
	\end{itemize}
	
	То есть
	\[
	(BA)(S)
	\]
	предкомпактно для любого ограниченного \(S\). Следовательно,
	\[
	BA\in\mathcal{K}(E).
	\]
\end{proof}

\begin{remark}[Линейность класса компактных операторов]
	Если
	\[
	A_1,A_2\in\mathcal{K}(E_1,E_2),
	\qquad
	\alpha,\beta\in\mathbb{K},
	\]
	то
	\[
	\alpha A_1+\beta A_2\in\mathcal{K}(E_1,E_2).
	\]
	Это следует из того, что сумма двух предкомпактных множеств предкомпактна, а умножение
	предкомпактного множества на число сохраняет предкомпактность.
\end{remark}

\begin{theorem}[Замкнутость класса компактных операторов по операторной норме]
	Пусть \(E_1\) --- нормированное пространство, \(E_2\) --- банахово пространство,
	\[
	A_n\in\mathcal{K}(E_1,E_2),
	\qquad
	A_n\to A
	\]
	по операторной норме:
	\[
	\|A_n-A\|\to0.
	\]
	Тогда
	\[
	A\in\mathcal{K}(E_1,E_2).
	\]
\end{theorem}

\begin{proof}
	По критерию через единичный шар достаточно доказать, что
	\[
	A\bigl(\overline{B}_1(0)\bigr)
	\]
	предкомпактно.
	
	Так как \(E_2\) банахово, достаточно доказать, что множество
	\[
	A\bigl(\overline{B}_1(0)\bigr)
	\]
	вполне ограничено.
	
	\textbf{1. Приблизим \(A\) компактным оператором.}
	
	Пусть
	\[
	\varepsilon>0.
	\]
	Так как
	\[
	\|A_n-A\|\to0,
	\]
	существует номер \(N\), такой что
	\[
	\|A_N-A\|<\frac{\varepsilon}{2}.
	\]
	Тогда для любого \(x\in\overline{B}_1(0)\)
	\[
	\|A_Nx-Ax\|
	\leqslant
	\|A_N-A\|\cdot\|x\|
	<
	\frac{\varepsilon}{2}.
	\]
	
	\textbf{2. Возьмём конечную сеть для \(A_N(\overline{B}_1(0))\).}
	
	Оператор \(A_N\) компактен, значит множество
	\[
	A_N\bigl(\overline{B}_1(0)\bigr)
	\]
	предкомпактно. Так как \(E_2\) банахово, оно вполне ограничено.
	
	Следовательно, существует конечная \(\varepsilon/2\)-сеть
	\[
	y_1,\ldots,y_m
	\]
	для множества
	\[
	A_N\bigl(\overline{B}_1(0)\bigr).
	\]
	То есть для любого \(x\in\overline{B}_1(0)\) найдётся номер \(j\in\{1,\ldots,m\}\), такой что
	\[
	\|A_Nx-y_j\|<\frac{\varepsilon}{2}.
	\]
	
	\textbf{3. Получим конечную сеть для \(A(\overline{B}_1(0))\).}
	
	Для такого \(x\) и соответствующего \(j\) имеем
	\[
	\|Ax-y_j\|
	\leqslant
	\|Ax-A_Nx\|+\|A_Nx-y_j\|
	<
	\varepsilon.
	\]
	Значит,
	\[
	y_1,\ldots,y_m
	\]
	является конечной \(\varepsilon\)-сетью для множества
	\[
	A\bigl(\overline{B}_1(0)\bigr).
	\]
	
	Так как \(\varepsilon>0\) произвольно, множество \(A(\overline{B}_1(0))\) вполне ограничено.
	Следовательно, оно предкомпактно.
	
	Значит,
	\[
	A\in\mathcal{K}(E_1,E_2).
	\]
\end{proof}

\begin{definition}[Вполне непрерывный оператор]
	Оператор \(A\in\mathcal{L}(E_1,E_2)\) называется \textbf{вполне непрерывным}, если из слабой
	сходимости
	\[
	x_n\weakconv x
	\]
	следует сильная сходимость образов:
	\[
	\|Ax_n-Ax\|\to0.
	\]
\end{definition}

\begin{theorem}[Компактный оператор улучшает слабую сходимость]
	Пусть \(E_1,E_2\) --- нормированные пространства и
	\[
	A\in\mathcal{K}(E_1,E_2).
	\]
	Если
	\[
	x_n\weakconv x
	\]
	в \(E_1\), то
	\[
	\|Ax_n-Ax\|\to0.
	\]
\end{theorem}

\begin{proof}
	Обозначим
	\[
	y_n=Ax_n.
	\]
	Нужно доказать, что
	\[
	y_n\to Ax
	\]
	по норме.
	
	\textbf{1. Последовательность образов предкомпактна.}
	
	Слабо сходящаяся последовательность ограничена. Поэтому множество
	\[
	\{x_n:n\in\mathbb{N}\}\cup\{x\}
	\]
	ограничено в \(E_1\).
	
	Так как \(A\) компактен, множество его образов
	\[
	\{Ax_n:n\in\mathbb{N}\}\cup\{Ax\}
	\]
	предкомпактно в \(E_2\).
	
	\textbf{2. Образы слабо сходятся к \(Ax\).}
	
	Так как \(A\) --- линейный ограниченный оператор, он сохраняет слабую сходимость:
	\[
	x_n\weakconv x
	\quad\Longrightarrow\quad
	Ax_n\weakconv Ax.
	\]
	То есть
	\[
	y_n\weakconv Ax.
	\]
	
	\textbf{3. Предкомпактность плюс слабая сходимость дают сильную сходимость.}
	
	Докажем от противного. Пусть
	\[
	y_n\not\to Ax
	\]
	по норме. Тогда существует \(\varepsilon_0>0\) и подпоследовательность \(\{y_{n_k}\}\), такая что
	\[
	\|y_{n_k}-Ax\|\geqslant\varepsilon_0
	\qquad
	\forall k.
	\]
	Но множество \(\{y_n\}\) предкомпактно, поэтому из \(\{y_{n_k}\}\) можно выделить
	подпоследовательность \(\{y_{n_{k_j}}\}\), сходящуюся по норме к некоторому \(y\in E_2\):
	\[
	y_{n_{k_j}}\to y.
	\]
	Сходимость по норме влечёт слабую сходимость:
	\[
	y_{n_{k_j}}\weakconv y.
	\]
	С другой стороны, так как вся последовательность \(\{y_n\}\) слабо сходится к \(Ax\), то
	\[
	y_{n_{k_j}}\weakconv Ax.
	\]
	Слабый предел единственен. Значит,
	\[
	y=Ax.
	\]
	То есть
	\[
	y_{n_{k_j}}\to Ax
	\]
	по норме, что противоречит условию
	\[
	\|y_{n_k}-Ax\|\geqslant\varepsilon_0.
	\]
	
	Следовательно,
	\[
	Ax_n\to Ax
	\]
	по норме.
\end{proof}

\begin{remark}[Обратное утверждение]
	Если пространство \(E_1\) рефлексивно, то компактность оператора эквивалентна полной
	непрерывности:
	\[
	A\in\mathcal{K}(E_1,E_2)
	\quad\Longleftrightarrow\quad
	x_n\weakconv x \Rightarrow \|Ax_n-Ax\|\to0.
	\]
	В общем случае обратное утверждение требует дополнительных условий.
\end{remark}
